\scanpage{21}
D'où $s_*(z^2)=t_*(\Theta^2c')$, où $t:A\times Y\mathbin{-}A$ est la projection. On a alors, du pt de vue degré
\[
\begin{aligned}
\unsure{\operatorname{Tr}|\operatorname{tr}}(z^2)&=\unsure{\operatorname{Tr}|\operatorname{tr}}(\Theta^2c')=\unsure{\operatorname{Tr}|\operatorname{tr}}(\Theta^2\cdot1_A\times c)\\
&=\unsure{\operatorname{Tr}|\operatorname{tr}}(p_*(\Theta^2)\cdot c)
\end{aligned}
\]
où $p$ est la projection $A\times Y\mathbin{-}Y$.

\margin{\diagram{\begin{tikzcd}[ampersand replacement=\&]
A\times Y\arrow[d,"t"']\\A
\end{tikzcd}}}

Nous allons calculer $p_*(\Theta^2)$, en décomposant $p$ en $p'\times\operatorname{id}_P$, où $p':A\times A\mathbin{-}A$ est la deuxième projection. On a
\[
\Theta=\Theta_0\times1_P\quad\text{d'où }\Theta^2=\Theta_0^2\times1_P\quad\text{d'où}
\]
\[
p_*(\Theta^2)=p'_*(\Theta_0^2)\times\operatorname{id}_{P*}(1_P)=p'_*(\Theta_0^2)\times1_P.
\]

\margin{\diagram{\begin{tikzcd}[ampersand replacement=\&]
A\times Y\arrow[r,"p"]\arrow[d,equal]\&Y\\
A\times\underbrace{A\times P}\arrow[r,dash,"{\unsure{-|\to}}" description]\&\underbrace{A\times P}
\end{tikzcd}}}

D'autre part
\[
\begin{aligned}
\Theta_0^2&=\deleted{\unsure{\Delta}}\Delta^2+0+0-2(0,0)-2(0,0)+2(0,0)\\
&=-2(0,0)
\end{aligned}
\]
compte tenu que $\Delta^2=0$ \inserted{(équi. alg.)} puisque $\chi(A)=0$.

d'où
\[
p_*(\Theta^2)=-2(0\times0\times P)
\]
et on trouve pr suite
\[
\unsure{\operatorname{Tr}|\operatorname{tr}}(s_*(z^2))=-2c(0,0)
\]
où on pose
\[
c(0,0)=\unsure{\operatorname{Tr}|\operatorname{tr}}(c\cdot(0\times0\times P))
\]
Il existe un $c$ tel que $c(0,0)=1$, p. ex. $c=A\times A\times(a)$, où $a\in P(k)$. Donc remplaçant \inserted{encore} $c$ pr $-c$, on \unsure{trouve} comme valeurs successivement $-2$, et $2$, qui ont des signes différents !!

\mspage{20}
